\begin{lemma}
Assume the hierarchy \(\noderef{ParameterHierarchy}\) for
\((\varepsilon,\varepsilon_1,\varepsilon_2,n_0)\).  Under the two-bite
probability law with
\[
M(n)=\noderef{TwoBiteNaturalReducedVertexCount}(n),\qquad
p=\frac12\sqrt{\frac{\log n}{n}},
\qquad
K(n)=\noderef{TwoBiteNaturalIndependenceScale}(1+\varepsilon,n),
\]
with high probability, for every \(K(n)\)-set \(I\),
\[
\left|\bigcup_{v\in S_I}\binom{X_v(I)}2\right|\le \varepsilon_1 K(n)^2.
\]
\end{lemma}
\begin{proof}
Unfolding \(\noderef{WithHighProbability}\), the goal is to prove that the
probability of the event in the statement tends to \(1\).  For each \(n\), put
\[
M_n=M(n),\qquad K_n=K(n),\qquad
p_n=\noderef{TwoBiteEdgeProbability}(\beta,n),
\]
and let \(\mathbb P_n\) denote the two-bite probability functional
\[
\mathbb P_n(E)=
\noderef{TwoBiteEventProbability}(n,M_n,p_n,E).
\]
This is exactly the probability law occurring in the Lean statement.  The
hypothesis \(\beta=\frac12\) identifies
\[
p_n=\noderef{TwoBiteEdgeProbability}\!\left(\frac12,n\right)
     =\frac12\sqrt{\frac{\log n}{n}},
\]
while the two rounding hypotheses identify
\[
M_n=\noderef{TwoBiteNaturalReducedVertexCount}(n),\qquad
K_n=\noderef{TwoBiteNaturalIndependenceScale}(1+\varepsilon,n)
\]
for every \(n\).
We will use the following elementary properties of this finite probability
law.  By unfolding \(\noderef{TwoBiteEventProbability}\), for any event \(A\),
\[
\mathbb P_n(A)=
\sum_{\omega}
\mathbf 1_A(\omega)\,\noderef{TwoBiteSampleWeight}(\omega).
\]
The weights are nonnegative, because
\(\noderef{TwoBiteSampleWeight}\) is the product of two
\(\noderef{GnpGraphWeight}\) factors and
\(\noderef{UniformInjectionWeight}\).  They have total mass \(1\), since the
red and blue \(G(M_n,p_n)\) product weights each sum to \(1\) and the
uniform-injection weights sum to \(1\) over the injection coordinate.  Hence
\[
\mathbb P_n(A^c)=1-\mathbb P_n(A),
\]
and the same indicator-sum formula gives monotonicity and subadditivity: if
\(A\subseteq B\), then \(\mathbb P_n(A)\le\mathbb P_n(B)\), and
\[
\mathbb P_n(A\cup B)\le \mathbb P_n(A)+\mathbb P_n(B).
\]

Let \(\mathcal S_n\) be the event
\[
\mathcal S_n(\omega):=
\noderef{TwoBiteSmallClosedPairsEvent}(k=K_n,\varepsilon,\varepsilon_1)(\omega).
\]
Unfolding \(\noderef{TwoBiteSmallClosedPairsEvent}\), this is
\[
\mathcal S_n(\omega)\quad\Longleftrightarrow\quad
\forall I\subseteq \mathrm{Fin}(n),\ |I|=K_n\Longrightarrow
\noderef{ClosedPairsBound}\!\left(
\left|\noderef{TwoBiteClosedPairsInClass}
  \bigl(\omega,I,\noderef{TwoBiteSmallClass}(\omega,\varepsilon,I)\bigr)
\right|,
\varepsilon_1,K_n\right).
\]
Unfolding \(\noderef{ClosedPairsBound}\), the final predicate is precisely
the displayed inequality
\[
\left|\bigcup_{v\in S_I}\binom{X_v(I)}2\right|
\le \varepsilon_1 K_n^2
\]
for that \(I\).  Thus \(\mathcal S_n\) is exactly the small closed-pairs event
claimed in the statement, and the desired conclusion is
\(\mathbb P_n(\mathcal S_n)\to1\).

Let
\[
\mathcal D_n(\omega):=\noderef{FiberAndDegreeEvent}(\omega,\varepsilon_2).
\]
The first conjunct of \(\noderef{ParameterHierarchy}\) gives
\(0<\varepsilon_2\).  Applying \(\noderef{FiberAndDegreeConcentration}\) with
this positivity proof, the same value of \(\beta\), the equality
\(\beta=\frac12\), and the same function \(M\) satisfying
\[
M(n)=\noderef{TwoBiteNaturalReducedVertexCount}(n),
\]
gives the high-probability statement
\[
\mathbb P_n(\mathcal D_n)\to1. \tag{1}
\]
No change of probability space is involved: the child theorem has the same
arguments \(n,M_n,\noderef{TwoBiteEdgeProbability}(\beta,n)\) as
\(\mathbb P_n\).  By the complement identity just recorded, (1) is
equivalently the statement that
\[
\mathbb P_n(\mathcal D_n^c)=1-\mathbb P_n(\mathcal D_n)\to0. \tag{1a}
\]

We next record the bad-event estimate needed for the high-probability
combination.  Fix \(n>n_0\).  Apply
\(\noderef{SmallClosedPairsBadEventProbability}\) with the same hierarchy
\(\noderef{ParameterHierarchy}(\varepsilon,\varepsilon_1,\varepsilon_2,n_0)\),
with the equality \(\beta=\frac12\), with this inequality \(n>n_0\), and with
the rounded equalities
\[
M_n=\noderef{TwoBiteNaturalReducedVertexCount}(n),\qquad
K_n=\noderef{TwoBiteNaturalIndependenceScale}(1+\varepsilon,n).
\]
The child theorem is stated for
\[
\noderef{TwoBiteEventProbability}
 \bigl(n,M_n,\noderef{TwoBiteEdgeProbability}(\beta,n),\cdot\bigr),
\]
which is exactly \(\mathbb P_n\) by the definition of \(M_n\) and \(p_n\).
The hypothesis \(\beta=\frac12\) is exactly the beta hypothesis required by
the child theorem and identifies its edge parameter with the rounded value
\(p_n=\noderef{TwoBiteEdgeProbability}(\frac12,n)\).
After rewriting \(k(n)\) to \(K_n\), the child event is exactly
\[
\{\omega:\mathcal D_n(\omega)\text{ and }\neg\mathcal S_n(\omega)\}
=\mathcal D_n\cap\mathcal S_n^c.
\]
Indeed, its first conjunct is
\(\noderef{FiberAndDegreeEvent}(\omega,\varepsilon_2)=\mathcal D_n(\omega)\),
and its second conjunct is the negation of
\[
\noderef{TwoBiteSmallClosedPairsEvent}
  (k=K_n,\varepsilon,\varepsilon_1)(\omega),
\]
which is precisely \(\neg\mathcal S_n(\omega)\) by the definition of
\(\mathcal S_n\).  Therefore the child theorem gives
\[
\mathbb P_n(\mathcal D_n\cap\mathcal S_n^c)
\le (n:\mathbb R)^{-1} \qquad (n>n_0). \tag{2}
\]
All finite-union details needed for this estimate are contained in
\(\noderef{SmallClosedPairsBadEventProbability}\): the exact decomposition
of \(\mathcal D_n\cap\mathcal S_n^c\) into fixed-set bad events, the empty
case \(K_n>n\), the \(\texttt{Nat.choose}(n,K_n)\) count, the finite union
bound using nonnegative \(\noderef{TwoBiteSampleWeight}\) terms, the
application of \(\noderef{SmallClosedPairsFixedSetTail}\), and the final
\(\noderef{ParameterHierarchyEventualComparisons}\) McDiarmid-union
comparison.

Now fix \(\rho>0\).  From \(\mathbb P_n(\mathcal D_n^c)\to0\), choose
\(N_{\mathcal D}\) such that, for all \(n\ge N_{\mathcal D}\),
\[
\mathbb P_n(\mathcal D_n^c)<\rho/2. \tag{3}
\]
The reciprocal bound is negligible in the elementary sense encoded by
\(\noderef{NegligibleProbability}\).  Explicitly, choose an integer
\(N_{\rho}\) with \(N_{\rho}>2/\rho\) and \(N_{\rho}\ge1\).  Then, for all
\(n\ge N_{\rho}\),
\[
0\le(n:\mathbb R)^{-1}\le (N_{\rho}:\mathbb R)^{-1}<\rho/2. \tag{4}
\]
Finally choose an integer threshold \(N\) with
\[
N\ge N_{\mathcal D},\qquad N\ge N_{\rho},\qquad N>n_0.
\]
Then every \(n\ge N\) satisfies \(n\ge N_{\mathcal D}\), \(n\ge N_{\rho}\),
and \(n>n_0\).  For such \(n\), the elementary set inclusion
\[
\mathcal S_n^c\subseteq \mathcal D_n^c\cup(\mathcal D_n\cap\mathcal S_n^c)
\]
holds: if \(\mathcal S_n\) fails, either \(\mathcal D_n\) fails, or
\(\mathcal D_n\) holds and \(\mathcal S_n\) fails.  Monotonicity and
subadditivity of the finite-sum probability \(\mathbb P_n\) apply to this
inclusion as follows:
\[
\mathbf 1_{\mathcal S_n^c}
\le
\mathbf 1_{\mathcal D_n^c}
+\mathbf 1_{\mathcal D_n\cap\mathcal S_n^c}.
\]
Multiplying by the nonnegative sample weights and summing gives
\[
\mathbb P_n(\mathcal S_n^c)
\le \mathbb P_n(\mathcal D_n^c)
  +\mathbb P_n(\mathcal D_n\cap\mathcal S_n^c).
\]
Using (3), (2), and (4), this gives
\[
\mathbb P_n(\mathcal S_n^c)
< \rho/2+\rho/2=\rho.
\]
Thus \(\mathbb P_n(\mathcal S_n^c)\to0\).  Using the complement identity for
the same probability law, this is equivalent to
\(\mathbb P_n(\mathcal S_n)\to1\).  This is exactly the
\(\noderef{WithHighProbability}\) assertion of the statement.
\end{proof}
